% ============================================================
%  chapter/examples.tex — Text patterns & citation showcase
%
%  A grab-bag of the devices books actually use: quotations, verse,
%  epigraphs, footnote-dense prose, and every common way to cite.
% ============================================================
\chapter{Text Patterns}\label{ch:patterns}

\section{Quotations}
Short quotations stay inline \enquote{like this one}; anything worth
its own lines uses the quotation environment:

\begin{quotation}
\noindent\itshape
Typography exists to honour content.\index{typography!purpose}
\end{quotation}
\begin{flushright}
--- \citeauthor{bringhurst2004}, \citeyear{bringhurst2004}
\end{flushright}

A displayed extract from another work keeps its original lineation
in the \texttt{verse} environment:

\begin{verse}
To see a World in a Grain of Sand\\
And a Heaven in a Wild Flower\\
\hspace*{2em}Hold Infinity in the palm of your hand\\
\hspace*{2em}And Eternity in an hour
\end{verse}

\section{Epigraphs}
Chapter openings often carry a short attributed line above the
first heading. There is no package for that here --- just a styled
block before the first section:

\begin{quote}
\small\itshape\color{AccentColor}
The form of a book must be answerable to the energy of its contents.
\end{quote}

\section{Citations, every common flavour}\label{sec:cites}
Biblatex offers one command per rhetorical situation, all reading
the same \texttt{reference.bib}. In running prose, name the source:
\textcite{knuth1984} designed TeX;
\textcite{lamport1994} put LaTeX on top of it. Parenthetical
asides take \parencite{mittal2005}; multiple sources collapse into
one bracket \parencite{knuth1984,lamport1994,mittal2005}.

When the year matters, write \textcite{texgyre}
(\citeyear{texgyre}); when the title itself is quoted,
\emph{\citetitle{knuth1984}} does the job. For a note at the bottom
of the page, use a footnote citation\footcite{bringhurst2004} ---
the full entry prints below the line without interrupting the
paragraph, a favourite in humanities and history
titles.\index{citations}

Every cited work lands automatically in the back-matter
bibliography, numbered in order of first appearance because the
template sets \texttt{sorting=none}.\index{bibliography}

\section{Footnotes in numbers}
Dense scholarly prose leans on notes.\footnote{Like this. Notes are
per-chapter here, resetting with each \texttt{\textbackslash chapter}.} They
should carry evidence,\footnote{Or the source, as footnotes do.}
qualification,\footnote{Or a polite disagreement.} or a joke too
good to cut\footnote{This one, arguably.} --- but never the argument
itself. If a reader must open your note to follow the sentence, the
sentence is in the wrong place.

\section{An index-rich close}
Everything a reader might hunt for belongs in the index, nested
where it helps:\index{book design}\index{book design!margins}%
\index{book design!grids}\index{page layout|seealso{book design}}
this closing section demonstrates sub-entries under a shared head,
and a \emph{see also} redirect. Compile, check
\texttt{main.ind}, and trust nothing you have not seen on a page.